\newcommandx{\supD}[3][1=P, 2=\setQ, 3=X]{\sup_{Q \in #2} \kl[Q][#1][#3]}
\newcommand{\argSupD}[1][\setQ]{\arg\min_{P}\sup_{Q \in #1}}
\newcommand{\myDiv}[1][P(X)]{\measure{D}{\mathrm{CD}}{\set{H} \,\middle\|\, #1}}
\newcommandx{\klExp}[3][1=Q, 2=P, 3=x]{\sum_{#3} #1(#3) \ln \frac{#1(#3)}{#2(#3)}}
\newcommand{\inSet}[1]{\mathFix{#1 \in \set{#1}}}

\newcommand{\saddle}{\mathFix{(\saddP, \saddR)}}
\newcommand{\fExp}{\sum_i R(i) \klExp[Q_i]}
\newcommand{\cDist}[1][]{\mathFix{\sub{C}{#1}}}
\newcommand{\ccDist}[1][]{\mathFix{\sub{C^{*}}{#1}}}
\newcommand{\cPrior}{\mathFix{c}}
\newcommand{\like}{Q}
\newcommand{\sadd}[1]{\mathFix{#1^{*}}}
\newcommand{\saddR}{\sadd{R}}
\newcommand{\saddP}{\sadd{P}}
\newcommand{\est}{T}

\newcommand{\fun}[1][]{\ifthenelse{\isempty{#1}}{}{(#1)}}

Consider a random variable $X$.
In probabilistic modeling, our knowledge about
$X$ is represented by a probability distribution, which serves as a basis for making inferences.

However, in many situations, our information may be too limited to justify choosing a single distribution.
Instead, we may only be able to specify a set of plausible distributions.
Such a set represents some form of \emph{uncertainty} or \emph{partial knowledge}: although it does not identify a
single distribution, it constrains the range of plausible probabilistic models.

A well-known example of this idea is the exchangeability assumption used in de Finetti’s theorem.
A sequence of random variables \infSeqExp{X} is said to be exchangeable, if the ordering of the variables carries no
information.
That is, the joint probability distribution remains invariant under permutations of the variables.
This assumption imposes a constraint, effectively introducing a set of plausible distributions.
Thus, exchangeability can be seen as a specific form of partial knowledge.

In this work, we adopt the view that a set of probability distributions can represent a meaningful state of partial
knowledge, and that such set-valued uncertainty may itself be used as the basis for inference.

More generally, partial knowledge may possess a hierarchical structure.
In the simplest case, uncertainty is represented by a single set of candidate distributions.
In richer settings, however, this uncertainty may be organized into nested collections, where elements of a set may
themselves be sets of distributions.
Such nesting can continue recursively, yielding a hierarchical uncertainty structure.

Crucially, the choice of hierarchical organization is itself informative: grouping distributions into subsets
expresses that certain alternatives are regarded as belonging to a common epistemic category and should therefore be
aggregated before being compared with alternatives in other groups.
Consequently, the hierarchy encodes information beyond the collection of leaf distributions alone.

This hierarchical organization may be interpreted as introducing multiple layers of latent structural decisions,
where higher-level choices restrict the admissible lower-level alternatives.

Within a Bayesian framework, inference ultimately requires a single probability measure.
Therefore, to use hierarchical partial knowledge in Bayesian inference, we require a principled method for
transforming a nested set-valued uncertainty representation into a single probability distribution while preserving
the informational structure encoded by the hierarchy.

We adopt the principle that the probability measure used to represent prior beliefs about the environment should
depend only on the available information about the environment itself, and not on agent specific considerations such as
the loss function or the informativeness of observations.
We refer to this modeling assumption as \emph{the principle of objectivity}.

From a minimax perspective, when uncertainty is represented by a set of candidate distributions, a natural
representative distribution is one that minimizes the worst-case divergence to the members of the set.
Such a distribution may be viewed as the most conservative or robust single-distribution summary of the available
partial knowledge.

\begin{definition}[Central Distribution]
    Let $X$ be a random variable and \setQ be a set of probability distributions over $X$.
    We say a distribution $C$ is a \emph{central distribution} of \setQ for $X$, if for every probability
    distribution $P$, we have
    \begin{equation}
        \supD[\cDist] \leq \supD[P],
        \label{eq:central-def}
    \end{equation}
    where $\measure{D}{\mathrm{KL}}{}$ denotes the Kullback–Leibler (KL) divergence.
\end{definition}

Note that, under this definition, the central distribution of \setQ is not necessarily an element of \setQ.
That is, the distribution that best represents the set in the specified sense may lie outside the set it summarizes.

Sometimes partial knowledge is represented by a hierarchical structure of sets of distributions.
Our current definition of the central distribution does not account for this case.
The following recursive definitions extend the concept to handle such situations.
While not directly about feature selection, this extension plays a key role in enabling the connection between
central distributions and feature selection, a connection that will become clearer later.

\begin{definition}[Hierarchy of Distributions]
    A \emph{hierarchy of distributions} over $X$ is defined recursively as follows:
    \begin{itemize}
        \item Any set of probability distributions over $X$ is a hierarchy of distributions over $X$.
        \item Any set whose elements are hierarchies of distributions over $X$ is itself a hierarchy of
        distributions over $X$.
    \end{itemize}
\end{definition}

Intuitively, a central distribution acts as a representative or summary of a set of distributions.
While it may not belong to the set itself, it reflects the overall behavior of the set in a way that is meaningful
for inference and also induces a notion of divergence for evaluating probability distributions under partial
knowledge.

The concept of a central distribution is not intended to yield a unique object; rather, all distributions satisfying
the definition can be regarded as equally good representatives.
However, when multiple central distributions exist, our constructions should not depend on an arbitrary
choice among them.
To avoid such dependence, we define the central divergence using the best divergence attained over the set of all
central distributions.

\begin{definition}[Central Divergence]
    \label{def:central-divergence}
    For any distribution $P$, we define the \emph{central divergence} of $P$ from a hierarchy of distributions \set{H},
    with respect to a random variable $X$, as follows:
    \begin{equation*}
        \myDiv = \inf_{\cDist} \kl[{\cDist}],
    \end{equation*}
    where $\cDist$ denotes a central distribution of \set{H} for $X$.
\end{definition}

So far, this definition applies only when the hierarchy is a single set of distributions.
The following definition extends the construction recursively to general hierarchical uncertainty structures.

\begin{definition}
    \label{def:hierarchy-of-distributions}
    Let $X$ be a random variable, and let \set{H^*} be a hierarchy of distributions over $X$ whose elements are
    hierarchies of distributions.
    We say \ccDist is a central distribution of \set{H^*}, if for every probability distribution $P$ we have:
    \begin{equation*}
        \sup_{\set{H} \in \set{H^*}} \myDiv[\ccDist(X)]  \leq \sup_{\set{H} \in \set{H^*}} \myDiv,
    \end{equation*}
    where $\myDiv[\none]$ refers to the central divergence as defined in Definition~\ref{def:central-divergence}.
\end{definition}

A simple example illustrates that hierarchical organization can affect the resulting representative distribution
even when the underlying admissible sets remain unchanged.

Let \(X\) be a random variable with three possible outcomes \(x_1,x_2,x_3\).
For each \(i\), let \(S_i\) denote the set of all probability distributions over \(X\) that assign zero probability to \(x_i\).

By symmetry, the central distribution of each set \(S_i\) is the uniform distribution over its support:
\[
    C_{S_1}=\left(0,\frac12,\frac12\right),\qquad
    C_{S_2}=\left(\frac12,0,\frac12\right),\qquad
    C_{S_3}=\left(\frac12,\frac12,0\right).
\]

If the three sets \(S_1,S_2,S_3\) are treated symmetrically as elements of a single set, then their central
distribution is
\[
    C_{\{S_1,S_2,S_3\}}=
    \left(\frac13,\frac13,\frac13\right).
\]

Now consider instead a hierarchical organization in which \(S_1\) and \(S_2\) are grouped together, while \(S_3\)
forms its own separate branch.

Taking the central distribution of these two representatives yields
\begin{align*}
    C_{\{\{S_1,S_2\}, S_3\}}
    &= \frac12 C_{\{S_1,S_2\}} + \frac12 C_{S_3} \\
    &= \frac14 C_{S_1} + \frac14 C_{S_2} + \frac12 C_{S_3} \\
    &= \left(\frac38,\frac38,\frac14\right).
\end{align*}

Thus, although the underlying admissible sets are identical in both constructions, the resulting representative
distribution differs.

This demonstrates that hierarchical structure carries informational content beyond the collection of leaf sets
alone.
Consequently, flattening a hierarchy into a single unstructured set may discard meaningful epistemic structure.

In many learning problems, uncertainty concerns not only the parameters of a predictive model but also its
structural form.
In particular, when the relevance of potential features is unknown, each candidate feature subset determines a
distinct set of distributions in which the target variable depends only on that feature subset.

Feature selection may therefore be viewed as uncertainty over a collection of distributions.
The hierarchical central distribution framework provides a natural mechanism for representing and aggregating such
uncertainty.
By organizing candidate features hierarchically, one may encode preferences over structural hypotheses, such as
favoring sparse subsets, grouping related attributes, or privileging particular interactions.

Candidate feature subsets typically induce overlapping families of admissible distributions.
Distributions exhibiting simpler dependency structure may belong simultaneously to many such families, whereas more
specific dependency patterns belong to fewer.
As a result, hierarchical aggregation through central distributions can induce a preference toward structurally
simpler models by assigning more weight to distributions compatible with multiple structural hypotheses.

With the notion of hierarchical central distributions established, we now develop the theoretical machinery needed
to characterize and compute them.
The following results provide optimization-theoretic characterizations of central distributions that form the
foundation of the subsequent analysis.

\begin{definition}[Interior Distribution]
    \label{def:interior}
    Let \setQ be a set of probability distributions over a random variable $X$.
    A distribution $R$ is called an \emph{interior distribution} of \setQ if, for every probability distribution $P$,
    there exists some $\inSet{Q}$ such that
    \begin{equation*}
        \kl[R] \leq \kl[Q],
    \end{equation*}
\end{definition}

\begin{theorem}
    Any distribution that can be expressed as a convex combination of elements of \setQ is an
    interior distribution of \setQ.
\end{theorem}

\begin{proof}
    Let $R = \sum_i \alpha_i Q_i$, where $Q_i \in \setQ$, $\alpha_i \geq 0$, and $\sum_i \alpha_i = 1$.

    Fix an arbitrary probability distribution $P$.
    By the convexity of KL divergence in its first argument, we have:
    \[
        \kl[R] \leq \sum_i \alpha_i \kl[Q_i] \cdot
    \]
    Now, note that the right-hand side is a weighted average of the KL divergences.
    So, there must be at least one index $j$ such that
    \[
        \kl[R] \leq \kl[Q_j] \cdot
    \]
    Since $Q_j \in \setQ$, the condition in Definition~\ref{def:interior} is satisfied,
    and $R$ is indeed an interior distribution of \setQ.
\end{proof}

\begin{theorem}
    Let $R$ be an interior distribution of \setQ.
    If $C$ is a central distribution of \setQ, it is also a central distribution of $\setQ \cup \{R\}$ and vice versa.
\end{theorem}

\begin{proof}
    This theorem follows directly from the definition of an interior distribution, which implies that for any
    probability distribution $P$, we have
    \[
        \supD = \supD[][\setQ \cup \{R\}] \cdot \qedhere
    \]
\end{proof}

One important consequence of this theorem is that, in order to identify a central distribution within a parametric
model, it suffices to restrict attention to the set of likelihood distributions induced by different values of the
model parameters.

The next theorem is a standard result that is typically used in the derivation of central distributions.

\begin{theorem}
    \label{th:estimator-to-var}
    Let $X$ and $T$ be two random variables, and let \setQ be a set of distributions over the pair $(X,T)$.
    Assume that for every \inSet{Q}, we have
    \[
        Q(X, T) = R(\XIT) Q(T),
    \]
    where $R$ is a fixed conditional distribution.
    If $C$ is a central distribution of \setQ for $T$, then $R(\XIT)C(T)$ is a central
    distribution of \setQ for $(X,T)$.
\end{theorem}

\begin{proof}
    Let $C$ be a central distribution of \setQ for $T$.
    For an arbitrary probability distribution $P$ we have
    \begin{equation}
        \label{eq:th5-proof}
        \sup_{Q \in \setQ} \Dkl{Q(X,T)}{R(\XIT) C(T)} \leq \sup_{Q \in \setQ} \Dkl{Q(X,T)}{R(\XIT) P(T)} \cdot
    \end{equation}

    Since the KL divergence is nonnegative, for every probability distribution $P$ and every \inSet{Q} we have
    \[
        \sum_t \Dkl{R(X \mid t)}{P(X \mid t)} Q(t) \geq 0.
    \]
    That implies
    \[
        \Dkl{Q(X,T)}{P(X,T)} - \Dkl{Q(X,T)}{R(\XIT) P(T)} \geq 0.
    \]
    Taking the supremum over \inSet{Q} preserves the inequality, since it holds pointwise for every \inSet{Q}:
    \[
        \sup_{Q \in \setQ} \Dkl{Q(X,T)}{R(\XIT) P(T)} \leq \sup_{Q \in \setQ} \Dkl{Q(X,T)}{P(X,T)} \cdot
    \]
    Combining this inequality with Inequality~\ref{eq:th5-proof}, we conclude that $R(\XIT)C(T)$ is a
    central distribution of \setQ for $(X,T)$.
\end{proof}


With the concept of hierarchical central distributions in place, we now develop the theoretical machinery needed
to characterize and compute them.
The following results provide optimization-theoretic characterizations of central distributions that form the
foundation of the subsequent analysis.

Throughout this section, all random variables are assumed to be discrete with finite support.
Before proceeding, we establish the existence of central distributions in the finite setting considered throughout
this section.

\begin{theorem}
    Let $X$ be a discrete random variable with finite support, and let \setQ be a finite set of probability
    distributions over $X$.
    Then a central distribution of \setQ for $X$ exists.
\end{theorem}

\begin{proof}
    For a fixed distribution $Q$, we may write
    \[
        \kl = c - \sum_x Q(x) \ln P(x).
    \]
    Since $\ln$ is upper semicontinuous on $[0,+\infty)$ and $X$ has finite support, it follows that
    $\kl$ is lower semicontinuous as a function of $P$.

    Because \setQ is finite, $\sFunc[P(X)][\setQ]$ is the maximum of finitely many lower semicontinuous functions
    and is therefore itself lower semicontinuous.

    The set of probability distributions over $X$ is a compact simplex.
    Hence, by the extreme value theorem, $\sFunc[][\setQ]$ attains its minimum.
    Therefore there exists a probability distribution \cDist such that
    \[
        \sFunc[\cDist(X)][\setQ] = \inf_P \sFunc[P(X)][\setQ].
    \]
    Thus $C$ is a central distribution of \setQ for $X$.
\end{proof}

The central distribution is essentially the solution to a convex--concave game.
Because solutions of convex--concave games correspond to saddle points,
by defining the appropriate objective function, we can characterize the central
distribution using standard convex optimization techniques.

\begin{definition} [Saddle Point]
    Consider a function $f:\set{w} \times \set{z} \to \field{R}$.
    We refer to a pair $(w^*, z^*)$ as a \emph{saddle point} for $f$, if we have
    \[
        f(w^*, z) \leq f(w^*, z^*) \leq f(w, z^*),
    \]
    for all \inSet{w} and \inSet{z}.
\end{definition}

Next, we establish the connection between saddle points and central distributions.
One significant insight provided by the next theorem is that it suggests inference and searching for a central
distribution can be combined and performed simultaneously when finding a saddle point.

\begin{theorem}
    % I've been unable to prove that the central distribution is a saddle point for f, or characterize conditions
    % that imply that.
    \label{th:saddle_is_central}
    Let $X$ be a discrete random variable, and \setQ be a finite set of distributions over $X$.
    Consider the following function:
    \begin{equation*}
        f(P, R)= \sum_i \kl[Q_i] R(i),
    \end{equation*}
    where $P$ and $R$ are probability distributions, $Q_i \in \setQ$ and the index $i$ ranges over all elements of \setQ.

    If \saddle is a saddle point for $f$, then \saddP is a central distribution of \setQ for $X$.
\end{theorem}

\begin{proof}
    Let \cDist be a central distribution of \setQ for $X$, and \saddle be a saddle point for $f$.
    Since $X$ is a discrete random variable and \setQ is finite, \cDist exists.
    As $C$ is a probability distribution, the point $(\cDist, \saddR)$ lies within the domain of $f$.
    By the definition of a saddle point, this implies that
    \begin{equation}
        \label{eq:th_proof_saddle_lq_dx}
        f\saddle \leq f(\cDist, \saddR).
    \end{equation}
    Because \saddR is a probability distribution, the definition of $f$ implies:
    \[
        f(\cDist, \saddR) \leq \sFunc[\cDist].
    \]

    On the other hand, by the definition of a saddle point
    \[
        f\saddle = \max_{R} f(\saddP, R),
    \]
    where $R$ is a probability distribution, and therefore
    \[
        \max_{R} f(\saddP, R) = \sFunc[\saddP].
    \]
    Thus, Equation~\ref{eq:th_proof_saddle_lq_dx} implies
    \[
        \sFunc[\saddP] \leq \sFunc[\cDist].
    \]
    Since \cDist is a central distribution, we must have \[\sFunc[\saddP] = \sFunc[\cDist],\] and $\saddP$ is also a
    central distribution of \setQ for $X$.
\end{proof}

\begin{lemma}
    \label{lm:saddle_conditions}
    Let $f$ be the function defined in Theorem~\ref{th:saddle_is_central}.
    The point \saddle is a saddle point for $f$, if and only if there are constants $\lambda$, $\eta$ and a
    function $\mu$, such that for every $i$, $x$ the following conditions hold:
    \begin{gather}
        \label{eq:kkt_lambda}
        \klExp[Q_i][\saddP][x] - \mu(i) = \lambda, \\
        \label{eq:kkt_w_avg}
        \saddP(x) = \sum_i \saddR(i) Q_i(x) \\
        \label{eq:kkt_r_sums_one}
        \sum_i \saddR(i) = 1, \\
        \label{eq:saddle_lm_r_geq_0}
        \saddR(i) \geq 0, \\
        \nonumber
        \mu(i) \geq 0, \\
        \label{eq:saddle_lm_rmu_is_zero}
        \mu(i) \saddR(i) = 0.
    \end{gather}

    Furthermore, for a fixed probability distribution $R$, a distribution $\hat{P}$ minimizes $f(\none,R)$, if and
    only if for all $x$ it satisfies
    \begin{equation}
        \label{eq:saddle_lm_minimizer_dist}
        \hat{P}(x) = \sum_i Q_i(x)R(i).
    \end{equation}
\end{lemma}

\begin{proof}
    Similar to $f$, we define $\hat{f}$ as follows:
    \begin{equation*}
        \hat{f}(P, R) = \fExp,
    \end{equation*}
    where the constraints that $P$ and $R$ be probability distributions are omitted.
    This enables us to formulate the saddle point of $f$ as the solution to a constrained optimization problem
    defined for $\hat{f}$.

    Let \saddle be a simultaneous solution to the following two optimization problems:

    \begin{equation}
        \begin{aligned}
            &\text{minimize} & &\hat{f}(P,R) & &\text{w.r.t.} \ P \\
            &\text{subject to} & &\sum_{x} P(x) = 1, \\
            \\
            &\text{maximize} & &\hat{f}(P,R) & &\text{w.r.t.} \ R \\
            &\text{subject to} & &R(i) \geq 0, & &\text{for all} \ i  \\
            & & &\sum_{i} R(i) = 1 &
        \end{aligned}\label{eq:twin_optimization}
    \end{equation}
    Clearly, if $\saddle$ exists, it will be a saddle point for $f$.

    Both of these optimization problems are convex and differentiable in the sense of convex
    optimization; see~\cite{boyd2004convex}.
    To establish this, consider that all the constraint functions are affine, and therefore are convex, concave and
    differentiable.

    Additionally, $\hat{f}$ is a differentiable function with respect to both $R$ and $P$.
    For any fixed $P$, $\hat{f}(P,\none)$ is an affine function and is concave.

    On the other hand, for a fixed solution to the second optimization problem denoted by $\hat{R}$, we can write
    \begin{equation}
        \label{eq:simplified_f}
        \hat{f}(P,\hat{R}) = c - \sum_{i,x} \hat{R}(i) Q_i(x) \ln P(x),
    \end{equation}
    where $c$ is not a function of $P$.
    Since $\hat{R}$ is a solution to the second optimization problem, it satisfies the corresponding constraints and
    is therefore a valid probability distribution.
    Hence, for all $i$, $x$ we have $\hat{R}(i)Q_i(x) \geq 0$, and $\hat{f}(\none, \hat{R})$  is
    the negative of a linear combination of concave $\ln P(x)$ functions with nonnegative coefficients.
    This function is convex.

    Thus, both optimization problems are convex and differentiable, with affine constraints.
    The constraints simply restrict the feasible set to the set of probability distributions, and it is always possible
    to find feasible points that meet all the constraints.
    Therefore, Slater’s condition is satisfied, and the KKT conditions provide necessary
    and sufficient conditions for optimality~\cite{boyd2004convex}.
    In other words, if there exists a point $\saddle$ that satisfies the KKT conditions for both problems, it will be a
    saddle point for $f$, and vice versa.

    The KKT conditions are as follows for every $x$, $i$ and constants $\lambda$, $\eta$:
    \begin{gather}
        \label{eq:kkt_proof_R_is_pxIe}
        \sum_i \saddR(i) Q_i(x) = \eta \saddP(x) \\
        \label{eq:kkt_proof_p_sums_one}
        \sum_{x} \saddP(x) = 1, \\
        \nonumber
        \klExp[Q_i][\saddP][x] - \mu(i) = \lambda, \\
        \label{eq:kkt_proof_r_sums_one}
        \sum_i \saddR(i) = 1, \\
        \nonumber
        \saddR(i) \geq 0, \\
        \nonumber
        \mu(i) \geq 0, \\
        \nonumber
        \mu(i) \saddR(i) = 0.
    \end{gather}

    By summing Equation~\ref{eq:kkt_proof_R_is_pxIe} over $x$ and using Equation~\ref{eq:kkt_proof_r_sums_one},
    we have
    \[
        \eta = 1\,.
    \]
    Thus, we get Equation~\ref{eq:kkt_w_avg} which also implicitly implies Equation~\ref{eq:kkt_proof_p_sums_one}.

    To prove the second part of the lemma,
    we observe that when $R$ is a probability distribution, Equations~\ref{eq:kkt_proof_R_is_pxIe}
    and~\ref{eq:kkt_proof_p_sums_one} are the KKT conditions for the minimization problem
    in Equation~\ref{eq:twin_optimization}.
    Since this minimization problem is convex and satisfies Slater's condition,
    the KKT conditions are necessary and sufficient.
    Hence
    \[
        \hat P(x)=\sum_i Q_i(x)R(i)
    \]
    if and only if $\hat P$ minimizes $f(\none,R)$, where $R$ is any fixed distribution.
\end{proof}

\begin{lemma}
    \label{lm:saddle-is-maximizer}
    Let \setQ be a finite set of probability distributions over a discrete random variable $X$.
    Define the function $g$ as:
    \[
        g(R) = \sum_i \kl[Q_i][R] R(i),
    \]
    where $Q_i \in \setQ$, and $R$ is a probability distribution such that \[R(X, I) = Q_i(X)R(I).\]
    Let $f$ be the function defined in Theorem~\ref{th:saddle_is_central}.
    If \saddle is a saddle point for $f$, then the joint distribution $Q_i(X)\saddR(I)$ is a maximizer of $g$.
\end{lemma}

\begin{proof}
    Let \saddle be a saddle point for $f$.
    By Lemma~\ref{lm:saddle_conditions}, \saddP satisfies Equation~\ref{eq:kkt_w_avg}.
    Letting $\saddR(X, I) = Q_i(X) \saddR(I)$ we have
    \[
        g(\saddR) = f\saddle.
    \]
    Now, let $R$ be an arbitrary probability distribution such that $g(R)$ is defined.
    From the second part of Lemma~\ref{lm:saddle_conditions}, it follows that $g(R)$ indeed corresponds to the
    minimum of $f(\none, R)$, which implies
    \[
        g(R) \leq f(\saddP,R).
    \]

    On the other hand, since \saddle is a saddle point, we also have
    \[
        f(\saddP, R) \leq f\saddle.
    \]
    Combining the two inequalities, we get:
    \[
        g(R) \leq g(\saddR).
    \]
    This shows that $g(\saddR)$ is the maximum of $g$, as desired.
\end{proof}

The following theorem gives a standard sufficient condition for optimality in the KL-minimax setting which is
the rationale behind the term ``central distribution.''
Like the center of a circle, the central distribution is equidistant from the boundary.

\begin{theorem} [Centrality]
    \label{th:centrality}
    Let $X$ be a discrete random variable, and let \setQ be a finite set of distributions over $X$.
    Suppose \cDist is a convex combination of elements of \setQ, with all mixture weights strictly positive.
    If there exists a constant $\lambda$ such that
    \begin{equation}
        \label{eq:equal-div}
        \kl[][\cDist] = \lambda,
    \end{equation}
    for every \inSet{Q}, then \cDist is the unique central distribution of \setQ for $X$.
\end{theorem}

\begin{proof}
    Lemma~\ref{lm:saddle_conditions} proves a set of sufficient conditions for a saddle point.
    Note that more restrictive conditions can also be imposed to derive an alternative set of sufficient conditions.
    Suppose additionally that $\saddR(i) > 0$ for every $i$.
    Then Equation~\ref{eq:saddle_lm_rmu_is_zero} implies $\mu(i) = 0$,
    and Lemma~\ref{lm:saddle_conditions} conditions can be simplified to
    \begin{gather*}
        \klExp[Q][\saddP][x] = \lambda, \\
        \saddP(x) = \sum_i \saddR(i) Q_i(x),
    \end{gather*}
    which must hold for every \inSet{Q} and some strictly positive probability distribution \saddR.
    Since \cDist is a convex combination of elements of \setQ, the corresponding convex
    weights define the distribution \saddR and \cDist satisfies these conditions.

    Therefore, by Lemma~\ref{lm:saddle_conditions} and Theorem~\ref{th:saddle_is_central}, \cDist is a central
    distribution of \setQ for $X$.

    On the other hand, for any probability distribution $C'$, we have
    \begin{align}\label{eq:cd-relation-to-lambda}
        \nonumber
        \kl[\cDist][C']
        &= \sum_i \saddR(i)\left(\kl[Q_i][C'] - \kl[Q_i][\cDist] \right) \\
        \nonumber
        &= \sum_i \saddR(i)\kl[Q_i][C'] - \lambda \\
        &\leq \sFunc[C'] - \lambda.
    \end{align}
    If $C'$ is a central distribution of \setQ for $X$,
    the definition of a central distribution implies
    \[
        \sFunc[C'] = \sFunc[\cDist] = \lambda.
    \]
    Therefore, by non-negativity of the KL divergence, we have
    \[
        \kl[\cDist][C'] = 0.
    \]
    Hence \cDist and $C'$ are identical probability distributions.
\end{proof}

\begin{corollary}
    \label{cor:C-maximizes-g}
    Under the assumptions of Theorem~\ref{th:centrality}, \cDist is a maximizer of the function $g$ defined in
    Lemma~\ref{lm:saddle-is-maximizer}.
\end{corollary}

\begin{proof}
    In the proof of Theorem~\ref{th:centrality}, we showed that $(\cDist, \saddR)$ is a saddle point for $f$.
    By letting $C(X, I) = Q_i(X)\saddR(I)$, Lemma~\ref{lm:saddle-is-maximizer} implies that \cDist is a maximizer
    of the function $g$.
\end{proof}

\begin{lemma}
    \label{lm:imperfection}
    Assuming the conditions of Theorem~\ref{th:centrality} and using the same $\lambda$,
    if $P$ is a convex combination of elements of \setQ, then there exists some $Q_0 \in \setQ$ such that
    \begin{equation*}
        \kl[Q_0][] \leq \lambda.
    \end{equation*}
\end{lemma}

\begin{proof}
    We proceed by contradiction.
    Suppose that for every \inSet{Q}, we have $\kl > \lambda$.
    This implies that
    \[
        g(P) > \lambda = g(C),
    \]
    where $C$ is the distribution specified in Theorem~\ref{th:centrality}, and $g$ is the function defined in
    Lemma~\ref{lm:saddle-is-maximizer}.
    Note that $g(P)$ is defined, because $P$ is a convex combination of elements of \setQ.

    However, by Corollary~\ref{cor:C-maximizes-g}, the distribution $C$ maximizes $g$, which contradicts the
    inequality $g(P) > g(C)$.
\end{proof}



Suppose we want to model an exchangeable sequence of random variables.
Let $S_n$ denote the sequence of the first $n$ variables: $S_n = (\seqExp{X})$.
As discussed earlier, exchangeability can be interpreted as a form of partial knowledge, which allows us to
represent our uncertainty using a central distribution over $S_n$.

However, this central distribution generally depends on the sequence length $n$.
According to the principle of objectivity, our beliefs should not depend on the specifics of our experimental setup,
and in particular, should not vary with the length of the sequence.

It is important to note that a distribution defined over a longer sequence induces marginal distributions over its
shorter subsequences.
This observation suggests that, in order to satisfy the principle of objectivity, we may instead determine a central
distribution over an infinite exchangeable sequence, from which the distributions of all finite initial segments can
be derived.

This requires constructing a central distribution that is well-defined over $S_{\infty}$.

\begin{definition}
    \label{def:infseq-central}
    Let \infSeq{S} be a sequence of random variables and \setQ be a set of probability measures,
    where each element of \setQ defines a distribution over each $S_n$.
    Assume \cDist[n] is a central distribution for $S_n$.
    Let \cDist be a probability measure that defines a distribution over every $S_n$.
    \cDist is a central distribution of \setQ for \infSeq{S}, if we have
    \begin{equation*}
        \toInf \kl[{\cDist[n]}][\cDist][S_n] = 0 \cdot
    \end{equation*}
\end{definition}

Intuitively, this definition states that as $n \to \infty$ the distribution that $C$ induces over $S_n$ becomes
increasingly similar to a central distribution for $S_n$.
In other words, \cDist is a central distribution for $S_{\infty}$.

Our next goal is to extend Theorem~\ref{th:estimator-to-var} so that it also applies in this asymptotic setting,
thereby providing a characterization of central distributions for infinite sequences.

\begin{theorem}
    \label{th:estimator-to-infseq}
    Let \infSeq{S} and \infSeq{T} be two sequences of random variables.
    Let \setQ be a set of probability measures, where each \inSet{Q} defines a joint distribution
    over $(S_n, T_n)$.
    Assume there exists a natural number $M$ and a fixed conditional distribution $R$
    such that, for all \inSet{Q} and $n \geq M$,
    \[
        Q(S_n, T_n) = R(S_n \mid T_n)Q(T_n).
    \]


    Let $C$ be a central distribution of \setQ for \infSeq{T} (Definition~\ref{def:infseq-central}).
    Define a probability measure $C^*$ by
    \[
        C^*(S_n, T_n) = R(S_n \mid T_n)C(T_n), \quad n \geq M.
    \]
    Then $C^*$ is a central distribution for \infSeq[]{(S_n, T_n)}.
\end{theorem}

\begin{proof}
    By Theorem~\ref{th:estimator-to-var}, for each $n \geq M$ we have
    \[
        C_n(S_n, T_n) = R(S_n \mid T_n) C_n(T_n),
    \]
    where $C_n(T_n)$ is a central distribution for $T_n$, and $C_n(S_n,  T_n)$ is a central distribution for
    $(S_n, T_n)$.
    Therefore,
    \begin{align*}
        \kl[{C_n}][C^*][S_n, T_n] &= \kl[{R(S_n \mid T_n)C_n}][R(S_n \mid T_n)C][T_n] \\
        &= \kl[{C_n}][C][T_n].
    \end{align*}
    Since $C$ is a central distribution for \infSeq{T} we have
    \[
        \toInf \kl[{C_n}][C][T_n] = 0.
    \]
    It follows that
    \[
        \toInf \kl[{C_n}][C^*][S_n, T_n] = 0,
    \]
    so $C^*$ is a central distribution for \infSeq[]{(S_n, T_n)}.
\end{proof}

We now present an asymptotic version of Theorem~\ref{th:lambda-centrality}, applicable to infinite sequences of
random variables.

\begin{theorem}
    \label{th:limiting-centrality}
    Let \infSeq{S} be a sequence of random variables.
    Suppose \cDist is a weighted average of elements of \setQ with strictly positive weights,
    and there exists a sequence of constants \infSeq{\lambda} such that
    \[
        \toInf  \sup_{\inSet{Q}} \left|\kl[][\cDist][S_n] - \lambda_n \right| = 0,
    \]
    then \cDist is a central distribution of \setQ for \infSeq{S}.
\end{theorem}

\begin{proof}
    For the purpose of this proof, we additionally assume that for each \(S_n\), there exists
    a central distribution $\cDist_n$ that satisfies the conditions of Theorem~\ref{th:lambda-centrality}.
    That is, there exists a constant \(\mu_n\) such that for every \inSet{Q}
    \[
        \kl[][\cDist_n][S_n] = \mu_n .
    \]

    For each \(n\), by the definition of central distribution, we have
    \begin{equation}
        \sup_{\inSet{Q}} \kl[][\cDist][S_n] \geq \mu_n,
        \label{eq:th_lim_central_proof_sup}
    \end{equation}
    and, by Lemma~\ref{lm:imperfection}, it also follows that
    \begin{equation}
        \inf_{\inSet{Q}} \kl[][\cDist][S_n] \leq \mu_n .
        \label{eq:th_lim_central_proof_inf}
    \end{equation}
    On the other hand, for each $n$ we have
    \[
        \left| \sup_{\inSet{Q}} \kl[][\cDist][S_n] - \lambda_n \right|
        \leq \sup_{\inSet{Q}}  \left|\kl[][\cDist][S_n] - \lambda_n \right|,
    \]
    and
    \[
        \left|\inf_{\inSet{Q}} \kl[][\cDist][S_n] - \lambda_n \right|
        \leq \sup_{\inSet{Q}}  \left|\kl[][\cDist][S_n] - \lambda_n \right|.
    \]
    This implies that the left hand sides of
    Equations~\ref{eq:th_lim_central_proof_sup} and~\ref{eq:th_lim_central_proof_inf}  both converge to $\lambda_n$.
    Now, by the squeeze theorem, the sequence \infSeq[]{\mu_n - \lambda_n} must converge to $0$.
    Therefore, for any arbitrary \(Q \in \setQ\), we have
    \begin{align*}
        \toInf \sum_{s_n} Q(s_n) \ln \frac{\cDist_n(s_n)}{\cDist(s_n)} &=
        \toInf (\kl[][\cDist][S_n] - \kl[][\cDist_n][S_n]) \\
        &= \toInf (\kl[][\cDist][S_n] -\mu_n + \lambda_n - \lambda_n) \\
        &= \toInf(\kl[][\cDist][S_n] - \lambda_n) - \toInf(\mu_n - \lambda_n) \\
        &= 0 .
    \end{align*}
    Since $\cDist_n$ satisfies the conditions of Theorem~\ref{th:lambda-centrality},
    $\kl[\cDist_n][\cDist][S_n]$ is a weighted average of
    $\sum_{s_n} Q(s_n) \ln \frac{\cDist_n(s_n)}{\cDist(s_n)} $, and we can conclude
    \[
        \toInf \kl[\cDist_n][\cDist][S_n] = 0 .
    \]
    Therefore, \cDist is a central distribution of \setQ for the sequence \infSeq{S}.
\end{proof}

\newcommand{\estLike}{Q \left( \est_{n} \mid \paramVal \right) }

Assume \setQ is a set of distributions that can be indexed by a parameter vector \param, and \infSeq{\est} is a
consistent estimator sequence for \param.
Let $P(\paramVal \mid\est_n = t)$ denote the posterior distribution obtained from an arbitrary smooth positive
prior.

Under regular parametric conditions, posterior distributions based on consistent estimators become asymptotically
insensitive to the choice of smooth positive prior.

Motivated by this phenomenon, we heuristically derive the asymptotic form of a central prior by considering the
posterior density evaluated at the point $t=\paramVal$.
A similar asymptotic construction also appears in~\cite{BERNARDO200517}.

Specifically, we define

\begin{equation}
    \label{eq:central-limit}
    \cPrior(\paramVal) \propto \toInf \frac{P(\paramVal \mid\est_n = t) \vert_{t = \paramVal}}{\mu_n},
\end{equation}
where \infSeq{\mu} is a sequence of positive constants chosen so that the limit exists for every value of \paramVal.

Let $k$ denote the normalizing constant converting proportionality into equality.
Since the logarithm is a continuous function, Equation~\ref{eq:central-limit} implies
\begin{equation*}
    \toInf \left(\ln \frac{P(\paramVal \mid\est_n = t) \vert_{t = \paramVal}}{\cPrior(\paramVal)}
               - \ln \mu_n + \ln k \right) = 0,
\end{equation*}

Since $\est_n$ is a consistent estimator, the sampling distribution $\estLike$ becomes increasingly concentrated
around $\est_n = \paramVal$ as $n \to \infty$.
Furthermore, under the usual regularity conditions, the posterior distribution becomes asymptotically
insensitive to the choice of smooth positive prior.
Motivated by this asymptotic prior-insensitivity, we heuristically replace the original prior with $c(\paramVal)$,
we obtain
\begin{equation*}
    \toInf \left( \sum_{t_n} Q(t_n \mid \paramVal) \ln
               \frac{Q(t_n \mid \paramVal)}{C(t_n)}
               - \ln \mu_n + \ln k \right) = 0 \cdot
\end{equation*}
Letting
\[
    \lambda_n = \ln \mu_n - \ln k,
\]
this suggests that, for each fixed \paramVal:
\begin{equation*}
    \toInf  \left( \Dkl{\estLike}{C(\est_n)} - \lambda_n \right) = 0 \cdot
\end{equation*}
Under additional regularity assumptions ensuring that this convergence is uniform over \paramVal, the assumptions of
Theorem~\ref{th:limiting-centrality} are satisfied.
Thus, the above asymptotic argument suggests that $\cPrior(\paramVal)$ acts as a central prior of \setQ for the
estimator sequence \infSeq{\est}.

A central distribution over an estimator cannot, by itself, be used to perform inference at the level of observed
outcomes.
To do so, we must convert it into a distribution over the primary variables of interest.
Theorems~\ref{th:estimator-to-var} and~\ref{th:estimator-to-infseq} provide a simple
method for this conversion.

\begin{example}[Binomial Experiment]
    \newcommand{\pv}{\paramVal}
    Consider a binomial experiment with $n$ Bernoulli trials.
    Denote the sequence of outcomes by $S_n$, and let $k$ be the number of successes observed in $S_n$

    Let \setQ be the set of distributions corresponding to different values of the success
    probability $\pv \in [0,1]$.
    We can consider $\pv$ as an index parameter for \setQ, and
    we write $Q(S_n \mid \pv)$ to denote the distribution $Q_{\pv}$ applied to $S_n$.
    Explicitly,
    \[
        Q(S_n \mid \pv) = \pv^k (1-\pv)^{n-k} \cdot
    \]

    Define $T_n$ to be the proportion of successes in $n$ trials.
    Note that $T_n$ is a consistent estimator for $\pv$.
    Therefore, we can use Equation~\ref{eq:central-limit} to find a central distribution for the sequence
    \infSeq{T}.

    Assuming a uniform prior on $\pv$, Bayes’ rule gives the posterior distribution
    \[
        P\left( \pv \, \middle| \, T_n = \tfrac{k}{n} \right) = (n+1)\binom{n}{k} \pv^k (1 - \pv)^{n - k},
    \]
    which is a $\mathrm{Beta}(k+1, n-k+1)$ distribution.
    Using Stirling’s formula, we derive its asymptotic form:
    \begin{equation*}
        P\left( \pv \, \middle| \, T_n = \tfrac{k}{n} \right) \sim
        \frac{(n+1)n^{(n+\hf)}}{\sqrt{2 \pi} (n-k)^{(n-k+\hf)} k^{(k+\hf)}} \, \pv^k (1 - \pv)^{n - k} \cdot
    \end{equation*}
    We now let $T_n \to \pv$:
    \begin{equation*}
        P\left(\pv \mid T_n = \pv \right) \sim \frac{n+1}{\sqrt{2 \pi n}} \, \pv^{-\hf}(1 - \pv)^{-\hf} \cdot
    \end{equation*}
    Therefore, if we define
    \[
        \mu_n = \frac{n+1}{\sqrt{2 \pi n}},
    \]
    the following limit exists:
    \[
        \toInf \frac{P\left(\pv \mid T_n = \pv \right)}{\mu_n} = \pv^{-\hf}(1 - \pv)^{-\hf} \cdot
    \]
    This implies that $c(\pv) \propto \pv^{-\hf}(1 - \pv)^{-\hf}$ is a central prior for \infSeq{T}.

    Now we apply Theorem~\ref{th:estimator-to-infseq} to obtain a central distribution for $\infSeq{S}$.
    Since $T_n$ is a deterministic function of $S_n$, the joint distribution $Q(S_n, T_n)$ coincides with the
    distribution of $S_n$: it assigns probability $Q(S_n)$ when $T_n$ is compatible with $S_n$, and zero otherwise.
    In other words, knowledge of $(S_n, T_n)$ is equivalent to knowledge of $S_n$.

    In the binomial experiment, all sequences with the same number of successes and failures have equal probability.
    Therefore,
    \[
        Q(S_n \mid \pv)
        = Q\!\left(S_n,\, T_n = \tfrac{k}{n} \mid \pv\right)
        = \frac{1}{\binom{n}{k}} Q(T_n \mid \pv) \cdot
    \]
    Let $C^*$ denote the central distribution of $\infSeq{S}$.
    By Theorem~\ref{th:estimator-to-infseq}, the distribution induced by $C^*$ on each $S_n$ is given by:
    \begin{equation*}
        C^*(S_n)
        %= C^*\!\left(S_n,\, T_n = \tfrac{k}{n}\right)
        = \frac{\gm{k+\hf} \gm{n-k+\hf}}{\gm{n+1}\gm{\hf}^2} \cdot
    \end{equation*}
\end{example}

When dealing with layered ignorance and attempting to determine the central distribution of a hierarchy of
distributions, as defined in Definition~\ref{def:hierarchy-of-distributions}, we typically encounter the case where the number of
elements at a given level of the hierarchy is finite.
In that setting, and under appropriate regularity conditions,
the central distribution reduces to an equally weighted average of the corresponding representative central
distributions.

\begin{example}[Fused Multinomial Model]
    \label{ex:fused-multinomial}
    \newcommand{\prodm}{\prod_{y=1}^m}
    \newcommand{\Qt}{Q_{\theta}}
    Let \infSeq{S} be a sequence where each $S_n = (\seqExp{X})$ represents $n$ independent trials of an
    experiment with $k$ possible outcomes,
    and \setQ be a set of probability distributions over $S_n$ indexed by a real vector \paramVal.
    Let $f: \set{X} \to \set{Y}$ be a deterministic function that maps each outcome of $X_i$
    to a label $Y_i = f(X_i)$, where $\set{Y} = \{1, \dots, m\}$.
    Assume that for all $i$ and $\Qt \in \setQ$, if $f(x_i) = f(x'_i)$ we have \[\Qt(\seqExp{x}) = \Qt(\seqExp{x'}).\]
    That is, the conditional distribution $\Qt(X_i \mid Y_i = y)$ is uniform over the preimage $f^{-1}(y)$.

    We can consider $\theta = (\theta_1, \dots, \theta_m)$ to be a probability vector over $\mathcal{Y}$,
    where each $\theta_y$ specifies the probability of observing label $y$.
    For a sequence $S_n$, let $n_y$ denote the number of times label $y \in \mathcal{Y}$ appears.
    Let $|f^{-1}(y)|$ denote the number of outcomes in $\mathcal{X}$ that are mapped to label $y$ by $f$.
    Then the probability of observing $S_n$ under $\Qt \in \setQ$, denoted $Q(S_n \mid \theta)$, is:
    \[
        Q(S_n \mid \theta) = \prodm \left( \frac{1}{|f^{-1}(y)|} \cdot \theta_y \right)^{n_y} \cdot
    \]

    Now, consider the estimator sequence $T_n = \frac{1}{n} \left( \seqExp[m]{n} \right)$, which is a consistent
    estimator for $\theta$.
    The conditional distribution $Q(S_n \mid T_n)$ is then given by:
    \[
        Q(S_n \mid T_n) = \frac{1}{n!} \prodm \frac{n_y!}{|f^{-1}(y)|^{n_y}} \cdot
    \]
    Note that this distribution is independent of $Q_\theta$.
    Therefore, by Theorem~\ref{th:estimator-to-var}, a central distribution for $S_n$ can be obtained from a central distribution
    for $T_n$.

    As $n \to \infty$, the distribution of $T_n$ converges to a multivariate Gaussian distribution.
    Using this fact and Equation~\ref{eq:central-limit}, we can determine a central prior for \infSeq{T}:
    \[
        c(\theta) \propto \prodm \theta_y^{-\hf}
    \]
    This is a Dirichlet prior and induces a Dirichlet-multinomial distribution for each $T_n$:
    \[
        C(T_n) = \frac{\gm{\hf[m]}\gm{n+1}}{\gm{n + \hf[m]}}
        \prodm \frac{\gm{n_y + \hf}}{\gm{\hf}\gm{n_y+1}} \cdot
    \]
    Note that here, $C$ denotes a central distribution for $T_\infty$, and $C(T_n)$ refers to the marginal
    distribution induced on $T_n$.

    By Theorem~\ref{th:estimator-to-var}, we can derive a central distribution for \infSeq{S}, which induces the
    following distribution over each $S_n$:
    \[
        C_f(S_n) = \frac{\gm{\hf[m]}}{\gm{n + \hf[m]}}
        \prodm \frac{\gm{n_y + \hf}}{\gm{\hf} |f^{-1}(y)|^{n_y}} \cdot
    \]
\end{example}
